\input{hand-preamble.tex}
\makeatletter
\newcommand{\hov}[1]{\mathord{\begin{tikzpicture}[baseline=(hv.base)]\node[inner sep=0pt,outer sep=0pt](hv){$#1$};
  \pgfmathsetmacro{\hvA}{(rnd-0.5)*1.2}\pgfmathsetmacro{\hvB}{(rnd-0.5)*1.2}%
  \draw[line width=0.6pt,line cap=round] ($(hv.north west)+(-0.5pt,1.6pt+\hvA pt)$) -- ($(hv.north east)+(0.8pt,1.6pt+\hvB pt)$);\end{tikzpicture}}}
\newcommand{\rt}[1]{\surd\!\hov{#1}}
\makeatother
\newcommand{\bul}{\par\hangindent=6mm\hangafter=1\noindent$\bullet$\ }
\newcommand{\nm}[1]{\par\hangindent=7mm\hangafter=1\noindent #1\ }
\begin{document}\shapka
\lk{10}{Проблемы оснований математики}
\term{Канторова теория множеств}\par
XIX в. --- работа по созданию оснований математики, МА (Б, Д, В)\par
XIX в. --- создание \term{теории множеств} (ТМ).\par
Г. Кантор (1845--1918) --- создатель ТМ.\par
Роль: ТМ --- универсальный фундамент математики.\par\vspace{1mm}
\red{Некоторые опр. и теоремы ТМ}
\bul Мн-во $E_1$ \term{равномощно} мн-ву $E_2$ ($|E_1|=|E_2|$), если сущ. биекция $E_1$ на $E_2$.
\bul $E_1$ \term{меньше или равно по мощности}, чем $E_2$ ($|E_1|\le|E_2|$), если сущ. инъекция $E_1$ в $E_2$.
\bul $E_1$ \term{меньше по мощности}, чем $E_2$ ($|E_1|<|E_2|$), если $|E_1|\le|E_2|$ и $|E_1|\neq|E_2|$.
\bul Мн-ва, равномощные $\mathbb{N}$, наз. \term{счётными}.
\bul Мн-ва, равномощные $\mathbb{R}$ (или $[0;1]$), наз. мн-вами \term{мощности континуума}.
\bul Th$_1$ (Кантор): $(E)(|E|<|P(E)|)$, где $P(E)$ --- \term{булеан} $E$.
\bul Th$_2$ (Кантор--Бернштейн): если $|E_1|\le|E_2|$ и $|E_2|\le|E_1|$, то $|E_1|=|E_2|$.
\bul Th$_3$ (монотонность): $E_1\subseteq E_2\Rightarrow|E_1|\le|E_2|$.
\par\vspace{1mm}
Но появились \term{парадоксы ТМ}, т.е. предъявляются утв., которые доказаны сами и их отрицания:
\nm{1.} \term{Парадокс Кантора} (1899 г.). Пусть $U$ --- мн-во всех мн-в. Тогда\par
\hspace*{6mm}1) По Th$_1$: $|U|<|P(U)|$, т.е. по опр. $|U|\neq|P(U)|$\par
\hspace*{6mm}2) $U\subseteq P(U)\Rightarrow$ по Th$_3$ $|U|\le|P(U)|$\par
\hspace*{14mm}$P(U)\subseteq U\Rightarrow$ по Th$_3$ $|P(U)|\le|U|$\ \ $\Rightarrow$ по Th$_2$: $|U|=|P(U)|$ --- против.
\nm{2.} \term{Парадокс Рассела} (1902 г.). Рассмотрим свойство мн-в $\beta(X)\leftrightarrow X\notin X$ (напр. $\mathbb{N}\notin\mathbb{N}$, $U\in U$) и мн-во $S=\{X\mid X\notin X\}$, т.е. $(X)(X\in S\leftrightarrow X\notin X)$. Таким образом $S\in S\leftrightarrow S\notin S$. Методом приведения к нелепости можно доказать, что $S\in S$ и $S\notin S$ --- противоречие.\par
\term{Парадокс брадобрея} --- популяризация парадокса Рассела.\par\vspace{1mm}
Так возникла \term{проблема непротиворечивости математики}: если „фундамент“ противоречив, то противоречива вся математика. Чтобы решить её, поставили 2 задачи:
\nm{1.} Задача доказательства непротиворечивости математики (но непонятно, что доказывать и как доказывать).
\nm{2.} Задача переосмысления средств логики и ТМ: их надёжности и универсальности.\par
Всё это ознаменовало \term{кризис оснований математики} (XIX--XX вв.):
\nm{1.} Открытие парадоксов.
\nm{2.} Расхождение математиков во взглядах на: 1) основные матем. понятия, принципы, логич. средства; 2) пути спасения математики.
\nm{3.} Принципиальная сложность обоснования непротиворечивости математики.\par\vspace{1mm}
\red{Критика некоторых принципов ТМ}
\nm{1.} \term{Принцип свёртывания} (ПС). Пусть $\rho(x)$ --- некот. свойство мн-в. Тогда сущ. мн-во, состоящее только из тех эл-тов, которые обладают свойством $\rho(x)$: $(\exists M)(x)(x\in M\leftrightarrow\rho(x))$.\par
ПС --- причина парадокса Рассела.\par
В аксиоматической ТМ Цермело--Френкеля ПС ограничен:\par
\hspace*{10mm}ОПС: $(A)(\exists M)(x)(x\in M\leftrightarrow\rho(x)\&x\in A)$.\par
При ОПС парадокс Рассела не возникает.
\nm{2.} \term{Абстракция актуальной бесконечности} (ААБ). Разрешается рассматривать бесконечные мн-ва как завершённые объекты вне зависимости от процесса их построения.
\nm{3.} \term{Аксиома выбора} (АВ). Для любого непустого семейства попарно непересекающихся непустых мн-в сущ. мн-во (\term{мн-во выбора}), содержащее ровно по 1 эл-ту каждого мн-ва.\par
\nb{Вместо ААБ используют \term{абстракцию потенциальной бесконечности} (АПБ): рассматривать бесконечные мн-ва как конечные, учитывая процесс их построения, и \term{абстракцию потенциальной осуществимости} (АПО) --- отвлечение от ограничений возможностей в материале, времени, пространстве.}\par\vspace{1mm}
\term{Доводы против АВ:}
\nm{1.} Как устроено мн-во выбора?
\nm{2.} „Парадокс“ Банаха--Тарского: трёхмерный шар равносоставлен двум своим копиям.\par
\term{Доводы в пользу АВ:}
\nm{1.} Очевидность.
\nm{2.} Частое использование в д-вах известных теорем: эквивалентность определений $\lim f(x)$.\par
Неизвестно, можно ли вывести АВ из других аксиом ТМ. \nb{(ответ в конце: АВ не зависит от остальных аксиом ZF)}\par\vspace{1mm}
\red{Критика некоторых логических средств}\par
Не обосновано использование в рассуждениях о бесконечных мн-вах следующих логич. средств: $A\vee\neg A$, $\neg\neg A\supset A$.\par
\term{Пример} (Дален). Докажем, используя $A\vee\neg A$, что сущ. иррац. числа $a$ и $b$, т.ч. $a^b$ --- рационально.
\nm{1.} $A$: пусть $(\rt{2})^{\rt{2}}$ --- рационально; тогда $a=b=\rt{2}$ --- иррац.
\nm{2.} $\neg A$: пусть $(\rt{2})^{\rt{2}}$ --- иррац.; тогда $a=(\rt{2})^{\rt{2}}$ и $b=\rt{2}$ --- иррац., но $a^b=2$ --- рац.\par
Не предъявили искомую пару чисел, т.е. не доказали конструктивно.\par
Док-во теорем вида $(\exists x)(P(x))$ наз. \term{эффективным}, если предъявлен объект (или алгоритм его построения), который удовлетворяет свойству $P(x)$, и \term{неэффективным} в ином случае.\par
Метод д-ва от противного --- источник неэффективных д-в теорем существования, т.к. для д-ва утв. $(\exists x)(P(x))$ допускается $\neg(\exists x)(P(x))$, получается противоречие, $\neg\neg(\exists x)P(x)$, и снимается двойное отрицание: $(\exists x)P(x)$. Но объект не предъявлен.\par\vspace{1mm}
\red{Программа Гильберта обоснования математики:}
\nm{1.} Доказать непротиворечивость математики.
\nm{2.} Сохранить ТМ Кантора.
\nm{3.} Найти контраргументы к критике.\par
Гильберт придерживался \term{финитной установки}: при д-ве непротиворечивости матем. нельзя использовать средства, которые подвергались критике в рассуждениях о бесконечных мн-вах.\par
Для д-ва непротиворечивости Гильберт использовал \term{метод формализации}: каждой аксиоматической неформальной теории $T$ сопоставить формальную теорию $T'$ с формальным языком и точно описанными логич. средствами. Гильберт хотел уточнить понятия „д-во“ и „теория“, а также превратить задачу д-ва непротиворечивости теории $T$ в матем. задачу: доказать, что не сущ. $T'$-формулы $F$, т.ч. $\vdash_{T'}F$ и $\vdash_{T'}\neg F$.\par\vspace{1mm}
\term{Судьба программы Гильберта}\par
В 1931 г. Курт Гёдель доказал, что непротиворечивость арифметики нельзя доказать средствами самой арифметики. Т.к. вся математика основывается на арифметике, это означает невыполнимость программы Гильберта в п.\,1.\par
Генцен в 1936 г. доказал непротиворечивость арифметики нефинитными средствами (построением дерева вывода).\par
Развиваются 2 логических типа мышления в математике:
\nm{1.} \term{Классический} --- использование ААБ, АВ, $A\vee\neg A$, $\neg\neg A\supset A$.
\nm{2.} \term{Интуиционистский} --- использование АПБ и АПО вместо ААБ, эффективное понимание утв. вида $A\vee B$ и $(\exists x)P(x)$, отказ от $A\vee\neg A$, $\neg\neg A\supset A$ как универсальных (в рассуждениях о бесконечном).\par
Также развивался \term{конструктивизм} (Марков, 1903--1979 гг.): те же логич. принципы, что и в интуиционизме, но также
\nm{1.} Рассмотрение только конструктивных объектов (слова; натуральные числа --- слова в алфавите палочек).
\nm{2.} Уточнение понятия эффективного алгоритма.\par
\end{document}
